\section{Sufficiency}\label{sec:sufficiency}
Sufficiency is a conjunction of floors on separately reported factors; it is recorded, never scored as a probability (Fig.~\ref{fig:suff}).
\begin{itemize}
\item \textbf{S0}: a label for ``an S1 floor fails''; it appears only inside \INS.
\item \textbf{S1}: at least one deterministic-confirmed \texttt{same} input valid on both versions; all changed units hit; no quarantined input that may hit a changed unit; no truncation, blind spot or tamper flag.
\item \textbf{S2}: S1 plus branch-direction coverage of changed regions of at least 0.80 per version, every contributing input deterministic-confirmed, and no lossy-equal observation on return or side channels.
\item \textbf{S3}: specified (plateau, trace stability, second explorer) but \textbf{not implemented}: it is refused in configuration and never emitted.
\end{itemize}
The 0.80 floor is an unvalidated default and no calibration was performed; the specification's miss-rate caps for S1 and S2 are claims to test, not results. Branch coverage counts conditional jumps, loop and send instructions only, with a conservative denominator. An independent review showed that S2 can be reached from two inputs, because levels do not depend on input count or diversity, and that the ``all declared channels instrumented'' condition is a constant true in the implementation. S1 and S2 must therefore not be read as a quality grade or as ``80\% of changed branch directions validated as behaviourally safe''. Sufficiency is also a contested proxy~\cite{inozemtseva2014coverage,just2014mut}.
\dotfig{sufficiency_levels}{Sufficiency levels S0--S2 (S3 not implemented)}{fig:suff}
